%% =====================================================================
%%  T-22-04  The Soothing Mirror
%%  -------------------------------------------------------------------
%%  One idea per file. Core is mandatory. Slots in canonical order.
%%  Max 700 words. No layout commands. No filler. New source material
%%  for this entry goes in sources/B3/T-22-04.tex -- never by
%%  rewriting this file (docs/RULES.md R0.1).
%%  =====================================================================
%% @kind: topic
%% @id: T-22-04
%% @title: The Soothing Mirror
%% @subtitle: Disarm and infuriate with the mirror effect --- reflect people until they meet themselves
%% @chapter: c22
%% @order: 40
%% @status: stable
%% @sources: B3
%% @updated: 2026-09-19

\topic{T-22-04}{The Soothing Mirror}{Disarm and infuriate with the mirror effect --- reflect people until they meet themselves}

\begin{core}
A mirror shows what is actually there --- which is what makes it such a clean instrument of deceit: mirror your enemies, doing exactly as they do, and they cannot decode your strategy. Held up to a psyche, the mirror seduces --- you share their values, you are their kind. Held up to behaviour, it teaches --- their conduct returned to them with interest. One instrument, two edges: the soothing reflection and the humiliating one.
\end{core}
\begin{mechanism}
B3's law builds the mirror as a family of moves. The seductive mirror: reflect the target's own values, taste and language back at them until the illusion of kinship does the persuading --- Marie Mancini's study of the young king, step one of every seduction, is a retreat into the target's own reflection. The shield-mirror: doing-as-they-do gives the opponent nothing to model --- you become unreadable because your moves are theirs, and the frustration of facing themselves makes them overreact. The teaching mirror: return mistreatment in exact measure, so the lesson arrives in their own handwriting and guilt does the rest --- what the source calls reflecting behaviour back until they feel guilty about it. And the mocking mirror, calibrated with what the source calls a touch of caricature: Ivan the Terrible enthroning the Tartar Simeon as a weak pretender to mirror the boyars' own childhood insult --- exaggeration added so the reflection cannot be missed. The authority behind it is the Shield of Perseus: face Medusa through the reflection, never directly. The warning is the source's own: beware mirrored situations --- two mirrors facing each other produce neither light nor image, only escalation.
\end{mechanism}
\begin{conditions}
Strongest in seduction, rivalry and discipline: wherever the target's own material is the most persuasive available --- their values, their style, their recent conduct. It weakens against the self-unaware (nothing to reflect), against the reflective few who recognise mirrors mid-move, and in intimate life, where mirroring shades into having no substance of your own, which reads as absence, then as mockery.
\end{conditions}
\begin{application}
  \item Study before you speak: seduction's first move is the retreat --- collect the target's vocabulary, heroes and grievances, and return them in their own accent.
  \item To disarm a rival, adopt their exact frame for one exchange. Nothing disorients like meeting themselves.
  \item To teach, return behaviour in kind and proportion --- once, visibly, without commentary. The lesson lands in their handwriting.
  \item Add caricature only when the plain mirror is being ignored. The exaggeration is spice, and it burns the dish if it leads.
  \item When you catch yourself mirroring everyone, stop and audit: a mirror with no face behind it is not a strategy, it is a symptom (T-28-01).
\end{application}
\begin{feedback}
  \item Working: the target warms fast and says \emph{it's like talking to myself}. The seductive mirror is holding.
  \item Working: the rival hesitates, visibly unable to model you. The shield is up.
  \item Failing: the reflection is read as mockery before it was meant as kinship. Calibration failed; the mirror is now a provocation.
  \item Failing: mirrored mistreatment escalates rather than teaches. You are inside a mirrored situation --- exit the hall, do not add angles.
\end{feedback}
\begin{failure}
The mirror practitioner's occupational disease is self-erasure: a person who wins every exchange by reflecting eventually has no unreflected position, and the values they borrowed become the ones they hold, second-hand and unexamined. The teaching mirror also has a failure loop: exactness of retaliation is hard to calibrate, and most mirrors over-return --- which converts the lesson into a feud with a receipt.
\end{failure}
\begin{countermeasures}
  \item When someone seems uncannily like you, count the echoes before you trust the kinship. The mirror builds its warmth fast and cheap (T-22-03).
  \item Against the shield-mirror: go concrete. Reflections cannot model specifics; ask for numbers, dates, and commitments and watch the frame crack.
  \item When your own conduct is returned to you once, visibly, without commentary --- take the lesson. It is the cheapest teaching you will ever receive.
  \item Refuse escalation in mirrored situations: the second retaliation is where feuds start. Break the symmetry or cap it, deliberately, in writing if needed.
\end{countermeasures}

\sources{B3}
\seesources
\seealso{T-22-01, T-07-04, T-21-06}
